\section{Introduction}
\label{sec:intro}

% Hook (PLAN.md section 2): paraphrase only; no quotation of the verse, no image of the character.
\gap{Hook paragraph.} The nickname of the character Chatur ``Silencer'' Ramalingam in the film \emph{3 Idiots}
\cite{hirani2009} rests on the belief that a silent release is the most offensive; in the film, a mock-Sanskrit verse
ranks releases by sound, from a loud blast down to the silent one, which it declares the deadliest.
An 1811 slang dictionary defines a related term as an escape ``more obvious to the nose than ears'' \cite{grose1811}.

\gap{Paragraph: the conjecture across cultures (32 languages, evidence levels; Figure~\ref{fig:culture}).}

\gap{Paragraph: three meanings of ``silent'' and three comparison bases:
(A) equal mass and composition, (B) as emitted, (C) among noticed events.}

\gap{Paragraph: contributions.}
(i) a mechanistic chain from sound to nose dose;
(ii) two transport theorems that bound what physics can do;
(iii) a selection model that produces the belief from a coin-flip truth;
(iv) a cross-cultural survey.
